\documentclass[11pt]{article}
\usepackage[margin=2.4cm]{geometry}
\usepackage[utf8]{inputenc}
\usepackage[T1]{fontenc}
\usepackage{amsmath,amssymb}
\usepackage{booktabs}
\usepackage{xcolor}
\usepackage{colortbl}
\usepackage{titlesec}
\usepackage{enumitem}
\usepackage{parskip}
\usepackage{fancyhdr}
\usepackage{tcolorbox}
\usepackage[colorlinks=true,linkcolor=blue!45!black,urlcolor=blue!45!black]{hyperref}

\definecolor{crit}{RGB}{155,25,25}
\definecolor{warn}{RGB}{180,110,0}
\definecolor{good}{RGB}{20,110,60}
\definecolor{hdr}{RGB}{30,55,105}

\titleformat{\section}{\normalfont\large\bfseries\color{hdr}}{\thesection}{0.6em}{}
\titleformat{\subsection}{\normalfont\normalsize\bfseries\color{hdr}}{\thesubsection}{0.6em}{}

\pagestyle{fancy}\fancyhf{}
\fancyhead[L]{\small\textit{Preregistro experimental}}
\fancyhead[R]{\small\thepage}
\renewcommand{\headrulewidth}{0.4pt}

\newcommand{\code}[1]{\texttt{\small #1}}
\newtcolorbox{cbox}[2]{colback=#1!5,colframe=#1,fonttitle=\bfseries,title={#2},
  boxrule=0.9pt,left=5pt,right=5pt,top=4pt,bottom=4pt}

\title{\textbf{\Large Preregistro experimental}\\[3pt]
\normalsize Qué pregunta queremos responder, y cómo sabríamos que la respondimos}
\author{\normalsize Repositorio \code{Ironsss/iag-symmetry}}
\date{\small 2026-08-27}

\begin{document}
\maketitle
\thispagestyle{empty}
\vspace{-1.2em}

% ============================================================
\section{El diagnóstico: qué faltaba}
% ============================================================

El proyecto ha producido un índice, una revisión crítica, un replanteamiento y un protocolo de auditoría. Lo que no ha producido nunca es una \textbf{pregunta enunciada de forma que las respuestas posibles sean enumerables}, con una hipótesis que pueda fallar, una operacionalización, controles, y un criterio de decisión declarado antes de medir.

El patrón ha sido consistente y es el problema de fondo: \emph{primero se construye un artefacto, después se busca qué significa}. Una metodología real invierte ese orden. Este documento lo invierte.

% ============================================================
\section{La evidencia de que esto importa: un confound producido hoy}
% ============================================================

Antes de escribir nada, se corrió el diagnóstico más barato posible: \emph{¿existe margen que estudiar en el régimen donde la asimetría per-flujo falla?} El resultado:

\begin{center}
\begin{tabular}{lr}
\toprule
Conjunto de features & AUC \\
\midrule
per-flujo asimetría (baseline) & 0.6863 \\
$+$ tamaño de grafo & 1.0000 \\
$+$ IAG (relacional) & 1.0000 \\
solo IAG & 1.0000 \\
\rowcolor{crit!12} solo tamaño de grafo (\code{n\_nodes}, \code{n\_edges}) & \textbf{1.0000} \\
\bottomrule
\end{tabular}
\end{center}

La última fila es el resultado. Dos enteros, sin ninguna relación con la asimetría, separan las clases perfectamente. La causa, verificada:

\begin{center}
\code{n\_nodes} en positivos: $51$ exacto \quad$\cdot$\quad en negativos: $50$ exacto \quad$\cdot$\quad \textbf{cero solapamiento}
\end{center}

Los positivos venían de corridas de infiltración y los negativos de corridas benignas \emph{separadas}. Las corridas de ataque tienen hosts atacantes extra. \code{n\_nodes} separa las \textbf{corridas}, no los \textbf{ataques}.

\begin{cbox}{crit}{Lo que casi se reporta}
Sin la fila de control ``solo tamaño'', la lectura natural de esa tabla habría sido: \emph{``el IAG alcanza AUC 1.0000 en el régimen sigiloso donde la asimetría per-flujo solo llega a 0.69''}. Un hallazgo espectacular, publicable, y completamente falso.

\medskip
El control estaba ahí por costumbre, no por diseño. Eso es precisamente la diferencia entre tener método y no tenerlo, y es la justificación de todo lo que sigue.
\end{cbox}

El diagnóstico, entonces, \textbf{no respondió la pregunta}. Demostró que el diseño era inválido. Ese es el punto de partida honesto.

% ============================================================
\section{La pregunta, derivada de lo medido}
% ============================================================

No se propone: se deriva de lo que el propio repositorio ya midió.

\textbf{Lo establecido.} La asimetría direccional per-flujo --- un solo escalar, $\log(1+\text{bytes}_\rightarrow) - \log(1+\text{bytes}_\leftarrow)$, sin grafo --- alcanza AUC $0.998$ en escaneo de puertos y $0.946$ en DDoS, y cae a $0.587$--$0.686$ en infiltración. Además, 24 de las 49 features ``clásicas'' de CICIDS ya son direccionales.

\textbf{El hueco.} Los ataques ruidosos ya los detecta cualquier cosa. Donde la detección importa --- y donde hay margen --- es el régimen sigiloso, y ahí la señal per-flujo está cerca del azar.

\begin{cbox}{hdr}{Pregunta}
\textbf{¿La asimetría medida a nivel \emph{relacional} aporta información de detección que las features direccionales \emph{per-flujo} no aportan ya, específicamente en el régimen donde las per-flujo fallan?}
\end{cbox}

Las respuestas posibles son enumerables: sí (con una magnitud), no, o ``no se puede saber con estos datos''. Las tres son resultados; la segunda es publicable; la tercera obliga a arreglar el instrumento.

% ============================================================
\section{Preregistro}
% ============================================================

Todo lo de esta sección se fija \textbf{antes} de correr el experimento. Un umbral elegido después de ver los números no es un criterio.

\subsection{Hipótesis y predicciones}

\textbf{H1.} En el estrato sigiloso (definido como aquel donde el baseline per-flujo obtiene AUC $\le 0.75$), agregar features relacionales produce una mejora $\Delta \text{AUC} \ge 0.05$ sobre el baseline.

\emph{Si H1 es cierta:} $\Delta \ge 0.05$, el intervalo de confianza excluye $0$, y la mejora \textbf{sobrevive a los cinco controles} de \S4.3.

\emph{Si H1 es falsa:} $\Delta < 0.05$, o el IC incluye $0$, o un control absorbe la mejora. Resultado negativo, informativo, y cierra la línea relacional.

\subsection{Operacionalización}

\begin{itemize}[leftmargin=1.3em,itemsep=1pt]
  \item \textbf{Unidad:} ventana temporal de 60\,s.
  \item \textbf{Etiqueta:} $y = 1$ si $\text{attack\_ratio} \ge \tau$, con $\tau = 0.01$ declarado, y sensibilidad reportada para $\tau \in \{0.005, 0.01, 0.05\}$.
  \item \textbf{Baseline:} estadísticos resumen del log-ratio direccional per-flujo dentro de la ventana (media, desv., p90, p99, fracción $>1$).
  \item \textbf{Tratamiento:} baseline $+$ features relacionales.
  \item \textbf{Métrica:} AUC y PR-AUC. IC 95\% bootstrap sobre la \emph{diferencia pareada}, no sobre cada AUC por separado.
\end{itemize}

\subsection{Controles --- cada uno descarta una explicación alternativa concreta}

\begin{center}
\begin{tabular}{clp{7.4cm}}
\toprule
& \textbf{Control} & \textbf{Explicación alternativa que descarta} \\
\midrule
C1 & Solo tamaño de grafo & ``El grafo es más grande'', no más asimétrico. \cellcolor{crit!8}\textbf{Es el que atrapó el confound de \S2.} \\
C2 & Diseño intra-corrida & Positivos y negativos de la \emph{misma} corrida, misma población de nodos. Descarta todo artefacto de configuración. \\
C3 & Per-flujo no-direccionales & ``Más features ayudan'', independientemente de que midan asimetría. \\
C4 & Ventanas degeneradas excluidas & Artefactos de borde (la ventana de 4 nodos y 2 aristas que ya invalidó H1 una vez). \\
C5 & Umbral de etiqueta declarado & Que el resultado dependa de dónde se puso $\tau$. \\
\bottomrule
\end{tabular}
\end{center}

\textbf{C2 es la corrección directa al error de \S2} y la más importante: mientras las clases provengan de corridas distintas, cualquier diferencia de configuración es un predictor perfecto.

\subsection{Regla de decisión}

\begin{enumerate}[leftmargin=1.5em,itemsep=1pt]
  \item $\Delta \ge 0.05$, IC excluye $0$, y la mejora de C1 es menor a la mitad de la del tratamiento $\Rightarrow$ \textbf{H1 sostenida}.
  \item $\Delta < 0.05$ o IC incluye $0$ $\Rightarrow$ \textbf{H1 rechazada}. La asimetría relacional no aporta aquí. Se reporta como negativo.
  \item C1 explica la mayor parte de la mejora $\Rightarrow$ \textbf{confounded}. No se reporta; se rediseña.
\end{enumerate}

\subsection{Potencia}

Medida por simulación sobre la diferencia de AUC, con baseline $0.69$:

\begin{center}
\begin{tabular}{rccc}
\toprule
$n$ por clase & $\Delta = 0.05$ & $\Delta = 0.10$ & $\Delta = 0.15$ \\
\midrule
50   & 77\% & 94\% & 100\% \\
100  & 85\% & 97\% & 100\% \\
\rowcolor{good!10} 250  & \textbf{95\%} & 100\% & 100\% \\
500  & 98\% & 100\% & 100\% \\
\bottomrule
\end{tabular}
\end{center}

\textbf{Requisito: $\ge 250$ ventanas por clase por estrato.} El experimento de \S2 tenía 60 por clase en test --- insuficiente incluso si el diseño hubiera sido válido. (Estimación de muestras independientes; una comparación pareada requiere menos, así que $250$ es conservador.)

% ============================================================
\section{¿Es ejecutable hoy? No --- y ahora se sabe exactamente por qué}
% ============================================================

El requisito de Fase 0 deja de ser ``arreglar el simulador'', que nunca estuvo especificado, y pasa a ser una lista derivada del diseño:

\begin{center}
\begin{tabular}{lcp{5.6cm}}
\toprule
\textbf{Requisito} & \textbf{Estado} & \textbf{Origen} \\
\midrule
Positivos y negativos intra-corrida & \cellcolor{crit!12}\textbf{violado} & C2 --- causa directa del confound de \S2 \\
Estrato sigiloso existente (AUC $0.6$--$0.75$) & \cellcolor{good!12}cumplido & medido: $0.686$ en infiltración \\
$\ge 250$ ventanas/clase/estrato & \cellcolor{crit!12}\textbf{violado} & potencia; hoy se obtienen $\sim 60$ \\
\code{attack\_ratio} cruzando $\tau$ & \cellcolor{crit!12}\textbf{violado} & mediana $0.926$; casi todo es $\sim 1$ o $0$ \\
Sin efecto techo en el baseline & \cellcolor{warn!12}parcial & falla en brute force (AUC $1.0$) \\
\bottomrule
\end{tabular}
\end{center}

\begin{cbox}{good}{Por qué esto es progreso y no otra postergación}
Las tres sesiones anteriores dijeron ``el simulador está roto'' sin poder decir \emph{en qué sentido} ni \emph{cuándo estaría arreglado}. Ahora hay cinco criterios de aceptación, cada uno derivado de un requisito del diseño experimental, cada uno verificable con una consulta.

\medskip
El generador ya no persigue realismo --- persigue estas cinco propiedades. Es una tarea de ingeniería acotada, no un problema de investigación abierto.
\end{cbox}

% ============================================================
\section{Dónde encaja el protocolo de auditoría}
% ============================================================

Queda claro su rol, que hasta ahora era ambiguo: \textbf{es el instrumento, no la contribución}.

Con 24 de 49 features direccionales correlacionadas entre sí, cualquier atribución ingenua sobre el resultado de este experimento daría la respuesta equivocada --- exactamente el colapso medido en \code{docs/protocolo/}. El protocolo existe porque la pregunta lo necesita para ser respondida, no al revés.

Si H1 se rechaza, el resultado publicable es doble: el negativo (la asimetría relacional no aporta sobre la per-flujo en el régimen que importa) más el instrumento que permitió establecerlo sin equivocarse.

% ============================================================
\section{Qué haría abandonar esta línea}
% ============================================================

Declarado ahora, para que no se decida después con los resultados a la vista:

\begin{itemize}[leftmargin=1.3em,itemsep=1pt]
  \item H1 rechazada con potencia adecuada ($n \ge 250$/clase) y controles limpios $\Rightarrow$ la línea relacional se cierra. Se reporta el negativo.
  \item El generador no alcanza los cinco criterios de \S5 en un esfuerzo acotado $\Rightarrow$ el bloqueante pasa a ser datos reales, declarado como único, y el proyecto se detiene ahí hasta conseguirlos.
  \item Un cuarto confound no anticipado sobrevive a los cinco controles $\Rightarrow$ el instrumento no es de fiar para esta pregunta y hay que cambiar de diseño, no de features.
\end{itemize}

\vspace{0.4em}
\begin{cbox}{hdr}{Estado}
Este documento no reporta ningún resultado sobre H1, porque el experimento no se ha corrido y no puede correrse todavía. Lo único que reporta es un confound propio, encontrado por un control, en un diagnóstico que no debía haberse construido así.

\medskip
Esa es la diferencia entre este documento y los tres anteriores: los otros presentaban hallazgos y buscaban su marco. Éste fija el marco y declara que no hay hallazgos aún.
\end{cbox}

\end{document}
